%%%%%%%%%%%%%%%%%%%%%%%%%%%%%%%%%
%
%       EXERCÍCIO
%
%%%%%%%%%%%%%%%%%%%%%%%%%%%%%%%%%

\definevalorquestao

\begin{exercicioBanco}[\valorquestao]
A temperatura média \(T\), em graus Celsius, de uma cidade ao longo de um dia é dada por
%
\[
T(h)=30-|14-h|,
\]
%
em que \(h\) representa a hora do dia, com \(0\le h\le 23\). Determine em que horário a temperatura atinge seu valor máximo.

% Define as alternativas
\newcommand{\alternativas}{%
    \begin{center}
        \begin{tabularx}{\textwidth}{XXXXX}
            (a) \(0\). &
            (b) \(6\). &
            (c) \(9\). &
            (d) \(14\). &
            (e) \(23\).
        \end{tabularx}
    \end{center}
}

% Define a resposta correta
\newcommand{\resposta}{D}

% Lógica condicional para exibição
\ifdefstring{\atividade}{lista}{%
    \alternativas
    \vspace{0.5em}
    
    \noindent\textbf{Resposta:} letra \textbf{\resposta}.
}{%
    \ifdefstring{\modo}{objetivo}{%
        \alternativas
    }{%
        % modo = discursiva → não mostra alternativas
    }
}
\end{exercicioBanco}

